% 第 22 章　磁场与运动电荷
\chapter{磁场与运动电荷}

\step{反常识开场}

\begin{mainline}
如果你的家里还留着一台老式的显像管电视机（或者在学校实验室里见过示波管），可以做一件让维修师傅皱眉的事：拿一块普通的冰箱贴磁铁，慢慢靠近屏幕。画面会立刻扭曲、变色，像有一只看不见的手把图像揉成了一团。把磁铁拿开，大部分扭曲会复原。

停下来想一想这件事有多奇怪。屏幕里的图像是电子打在荧光粉上画出来的。那块磁铁没有碰到任何一颗电子——它和电子之间隔着空气、隔着玻璃。更奇怪的是，电子明明在以每小时几千万公里的速度狂奔，磁铁却没有让它们变慢，也没有让它们变快，只是把它们的飞行路线\textbf{掰弯}了。

还有一件更挑剔的事。把一粒带电的小球（用毛衣摩擦过的气球就行）静悄悄地放在强磁铁旁边：什么也不会发生，磁铁对它视而不见。可一旦让这粒电荷动起来，哪怕只是缓慢飘过，磁场立刻“出手”了。磁铁仿佛能分辨电荷\textbf{动不动}，只对运动着的电荷感兴趣。

而最反直觉的是出手的方向。你推着一辆购物车往前走，想让车加速，就得顺着速度的方向推；这是我们全身的肌肉记忆。磁场偏偏不：电荷向北飞，它偏要向东推；电荷向东飞，它偏要向南或向北推。总之，它永远横在速度的路线上，像一根只拦路、不助跑的栏杆。世界上怎么会存在这样一种“永远不肯顺着速度使劲”的力？它拦来拦去，到底图什么？

再往大了看。地球两极的夜空中，极光像绿色的窗帘一样垂挂、摆动。极光里的光，是几十上百公里高空的大气原子被从太阳飞来的带电粒子撞出来的。可太阳风明明是朝整个地球吹的，为什么这些粒子偏偏挤到南北两极附近才掉下来？是谁在半路上给它们指了路？

这一章，我们就来认识这位只拦路、不助跑的选手：磁场。
\end{mainline}

\step{先猜一猜}

\begin{mainline}
往下读之前，先押三注。不用算，凭直觉给每句话判“对”或“错”，写下来：

\begin{itemize}
  \item 第一题：把一个带正电的小球用丝线吊着，静止地挂在一块强磁铁的北极附近（小球和磁铁之间没有接触），小球会被吸过去或者推开。
  \item 第二题：一个带正电的粒子，沿着指南针指针所指的方向（也就是磁感线方向）笔直射出去，它受到的磁力最大。
  \item 第三题：把带电粒子关在一个很强的匀强磁场里，让它转上一百万圈，只要磁场足够强、时间足够长，粒子的速率（快慢）总会被磁场越推越快。
\end{itemize}

剧透一句：三句话全是错的。第一题错在“静止”——磁场对不动的电荷完全无感；第二题错在“沿着”——磁力恰恰在最不顺的方向上才最大；第三题错得最值钱，它藏着一个深刻的分工：\textbf{磁场从不改变速度的大小，只改变速度的方向}。第三题还有一个推论值得你提前咀嚼：如果磁场永远不能把粒子推快，那么医院里回旋加速器不断提速的电量是谁出的？答案预告：另有其人，磁场只负责让粒子“留下来继续挨推”。这三题错在哪里，就是本章的主线。读完请回来自己批改。
\end{mainline}

\step{动手实验}

\begin{experiment}[让电流给指南针“下命令”]
你需要：一节五号电池（或充电宝加一根旧 USB 线），一米左右的导线，一个指南针——手机上指南针应用也行，但请先远离音箱和磁铁校准好。

第一步，立一个基准。把指南针放在桌上，等指针稳定指向南北。把导线拉直，沿着南北方向架在指南针正上方几厘米处（可以用两本书把导线垫起来）。

第二步，短暂通电。把导线两端接到电池正负极上，\textbf{接通两三秒就立刻断开}——这几乎是短路，导线会发热，千万别一直接着。通电的那两三秒里盯紧指针：它会猛地偏转一个角度；断电后它慢慢晃回南北。再把电池正负极对调，重复一次，指针会偏向另一边。

第三步，改变几何关系。把导线转成东西方向再通电，偏转明显变小甚至看不出来。把导线从指南针上方挪到正下方，偏转方向也会反过来。

这个实验就是 1820 年奥斯特在讲台上偶然看到的现象：\textbf{电流能让磁针偏转}。注意三个细节——不通电就没有偏转，说明动作的不是电池“漏电”到针上；电流反向偏转反向，说明这个作用有明确的方向性；导线方向与针的相对摆法影响效果，说明它强烈依赖几何关系。这三条线索，是我们发明“磁场”概念的全部起点。

安全提示：只能用干电池这类低电压电源，并且每次通电不超过几秒；绝对不要把导线插进墙上的插座。
\end{experiment}

\step{旧模型遇到麻烦}

\begin{mainline}
到目前为止，这本书给你的电磁学工具箱里只有一件工具：电荷与电荷之间的静电力——同性相斥、异性相吸，力沿着两者的连线。拿它来解释指南针实验，会连碰三鼻子灰。

麻烦一：静电力不挑动与不动。库仑定律里只有电荷量和距离，没有速度这一项。一个带电小球不管挂在磁铁旁边纹丝不动还是高速掠过，静电力理应一模一样。可实验事实是：静止时磁场不理它，一动起来就挨推。旧模型里根本没有“速度”这个位置可以安放这条新规律。

麻烦二：力的方向对不上。两个带电体之间的静电力永远沿着它们的连线——要么拉近，要么推远。可磁场推运动电荷的方向，既不从磁铁指向电荷，也不从电荷指向磁铁，而是垂直于电荷的运动方向，是一种“横着来”的力。连线方向的框架装不下它。

麻烦三，也是最致命的一击：两根通电线之间的力。取两根平行的导线，都通上同方向的电流，它们会互相吸引；电流一反一正，就互相排斥。请你用静电模型算算这笔账：导线是铜做的，铜里带正电的原子核和带负电的电子数量精确相等，整根导线是电中性的。两根中性的导线之间，静电力应该是零——不多不少的零。可它们偏偏相吸了。这个力不算小：工厂里意外短路的粗电缆会被这股力猛地抽打、甩弯，配电柜的铜排必须用坚固的绝缘支架牢牢锁住，就是怕短路电流把铜排生生掰变形。

三笔账都算下来，结论只有一个：电中性的导线因为\textbf{电流}而相互作用，电流就是“运动着的电荷”。运动电荷之间多出来一种全新的相互作用，它不在静电学的账本里。旧模型不是算错了一个数，而是少了一个玩家。
\end{mainline}

\step{发明新概念}

\begin{mainline}
缺一个玩家，我们就造一个。历史走的就是这条路：既然运动的电荷周围多出一种作用，那就说——\textbf{电流和磁铁在周围的空间里制造了一种“场”，叫磁场}；另一个运动电荷走进这片空间，就会被场推一把。

先做一件记账的工作：给磁场画地图。拿一枚小指南针，放在磁铁周围的任意一点，等指针稳定下来，记下它北极的指向；换个位置再记。把所有这些指向连成光滑的曲线，就得到\textbf{磁感线}：线上每一点的切线方向，就是该处磁场的方向；线画得越密，表示磁场越强。在磁铁外部，磁感线从北极出发、绕回南极；它们不像电场线那样起于正电荷、终于负电荷——事实上没有人找到过一个单独的“磁荷”，把磁铁掰成两半，你得到的不是单独的南极和北极，而是两块各有南北极的小磁铁。磁感线永远是闭合的环线，没有起点和终点。

这里值得请出一个比喻。磁感线像田径场四周观众举着的指示箭头，告诉你“朝这边走”——每一点的箭头给出方向。但它不像的地方更要紧：磁感线不是挂在空间里的真实丝线，空间里并没有一排排小箭头；它只是我们记录“此处磁场朝哪、多强”的记账方式，就像地图上的等高线不是山上画出来的白线。

方向地图有了，接下来是核心问题：磁场怎么推一个运动电荷？把大量实验结果压缩成三条：

\begin{itemize}
  \item \textbf{静止不推}。电荷不动，磁力为零。磁场只对运动电荷出手。
  \item \textbf{横着推}。磁力的方向永远同时垂直于电荷的速度方向和磁场方向。电荷顺着磁感线方向飞，磁力也是零；飞得越“横”，推得越狠。
  \item \textbf{认正负}。负电荷受的力与正电荷完全相反，像磁场对它们俩用了左右两只手。
\end{itemize}

这三条合成一句话：磁场对运动电荷施加一个\textbf{垂直于运动的侧向推力}。它有个名字，叫洛伦兹力。

洛伦兹力不直接推导线，但导线里跑着亿万个运动电荷。当电流通过放在磁场里的导线时，每个定向移动的电荷都挨一记横推，它们把这一记记推撞传递给整个晶格，宏观上看，就是\textbf{整根通电导线受到一个侧向的力}。这就是奥斯特实验的下半场：电流能推指南针，磁铁也能推电流——你的电风扇、电动牙刷、电动车的电机，全部靠这一下横推转起来。

把这句话变成一台会转的机器，只差一个线圈。把导线绕成矩形线框，架在两块磁铁的磁场中间，让电流沿框流动：线框的左右两条边里电流方向相反，受的磁力也就相反——一边被向上推，另一边被向下推，两个力一前一后拧着线框，线框就转起来了。麻烦在于：转到半圈之后，原来朝上的那条边转到了下面，受力还是朝上，合力开始往回拧，线框只会在平衡位置附近来回晃，成不了电机。 十九世纪的工程师用一个简单到近乎狡猾的零件解决了它——\textbf{换向器}：把线圈的两端接在两个半圆形的铜环上，电刷每转过半圈就自动把电流方向翻转一次。电流一翻，受力跟着翻，往回拧的力瞬间变成继续往前拧的力。于是线框被一股“永远朝同一个转向拧”的力矩驱着不停转下去。拆开任何一个直流小电机（玩具车里的那种），你都能看到这三件套：磁铁、线圈、换向器。扬声器是同一个原理的直线版：音圈套在永磁铁的磁场里，电流随音乐信号忽大忽小、忽正忽反，音圈就被推着前后振动，带动纸盆把振动交给空气。

还差最后一块拼图：指南针自己为什么指北？把一根通电导线绕成一个小线圈，放进磁场，线圈会转动，直到它的“面”朝向某个特定方向才稳定——它表现得和一枚磁针一模一样。于是可以给小线圈定义一个\textbf{磁矩}：方向由电流绕行方向决定，大小正比于电流乘以线圈面积。磁针就是一个小小的磁矩，地球本身是一个巨大的磁矩——地磁场在地面附近大约只有五万分之一特斯拉，弱得可怜，却足以让一枚几乎无摩擦的指针乖乖排队。物质磁性的微观线索也藏在这里：原子里的电子本身就带着微小的磁矩。注意，这不是说电子像行星绕太阳那样沿确定轨道转圈形成“小电流环”——量子力学告诉我们电子没有那种确定路线；电子的磁矩主要是它\textbf{内禀}的性质（叫自旋磁矩），就像它天生带电一样天生带磁。铁之所以特殊，是因为它的原子磁矩能在一大块区域里自发地整齐排队；普通材料里这些磁矩东倒西歪，互相抵消。
\end{mainline}

\begin{figure}[htbp]
\centering
\begin{tikzpicture}[scale=1.0]
  % 磁铁两极
  \draw[thick,fill=blue!8] (-4.2,-1.6) rectangle (-3.4,1.6);
  \draw[thick,fill=red!8] (3.4,-1.6) rectangle (4.2,1.6);
  \node at (-3.8,0) {N};
  \node at (3.8,0) {S};
  % 磁场线
  \foreach \y in {-1.0,0,1.0}
    \draw[-{Stealth[length=2.5mm]},blue!60!black] (-3.3,\y) -- (3.3,\y);
  \node[blue!60!black,anchor=south] at (0,1.15) {$\bm B$};
  % 线圈（斜放的矩形，透视画法）
  \draw[very thick,brown!75!black] (-1.3,-1.2) rectangle (1.3,1.2);
  % 电流方向
  \draw[-{Stealth[length=2.5mm]},thick,green!45!black] (-1.3,0.5) -- (-1.3,-0.5);
  \draw[-{Stealth[length=2.5mm]},thick,green!45!black] (1.3,-0.5) -- (1.3,0.5);
  \node[green!45!black,anchor=east] at (-1.45,0) {$I$};
  % 受力
  \draw[-{Stealth[length=3mm]},very thick,orange!85!black] (-1.3,0) -- (-1.3,1.9)
    node[left] {$\bm F$};
  \draw[-{Stealth[length=3mm]},very thick,orange!85!black] (1.3,0) -- (1.3,-1.9)
    node[right] {$\bm F$};
  % 转轴与转向
  \draw[dashed] (0,-2.2) -- (0,2.2);
  \draw[-{Stealth[length=2.5mm]},thick] (0.9,1.75) arc (60:120:1.6);
  \node[anchor=west] at (4.5,-1.9) {两边电流相反、受力相反，};
  \node[anchor=west] at (4.5,-2.5) {合力拧着线圈绕虚线轴转动。};
\end{tikzpicture}
\caption{电动机的核心：磁场中的通电线圈，两条竖边受方向相反的磁力，被拧着转动；换向器每半圈翻转电流，让拧劲始终朝同一转向。}
\end{figure}

\begin{figure}[htbp]
\centering
\begin{tikzpicture}[scale=1.05]
  % 磁场：指向纸内
  \foreach \x in {0,1.4,2.8,4.2,5.6}
    \foreach \y in {0,1.4,2.8}
      \node[blue!60!black] at (\x,\y) {$\otimes$};
  \node[blue!60!black,anchor=west] at (6.1,2.8) {$\bm B$：磁场，垂直指向纸内};
  % 电荷与速度
  \fill[red!70!black] (2.8,1.4) circle (3pt);
  \node[red!70!black,anchor=north] at (2.8,1.15) {$+q$};
  \draw[-{Stealth[length=3mm]},very thick,green!45!black] (2.8,1.4) -- (4.6,1.4)
    node[midway,below] {$\bm v$};
  \draw[-{Stealth[length=3mm]},very thick,orange!85!black] (2.8,1.4) -- (2.8,3.2)
    node[midway,left] {$\bm F$};
\end{tikzpicture}
\caption{洛伦兹力的方向：正电荷在指向纸内的磁场中向右运动，受力向上。三者两两垂直。若换成负电荷，$\bm F$ 掉头向下。}
\end{figure}

\step{一句话模型}

\begin{oneline}
磁场是一位只拦路、不助跑的裁判：它只推“动着的”电荷，而且永远横着推——推得越狠，电荷转弯越急，可它的快慢一丝一毫都不会变。
\end{oneline}

\begin{mainline}
这句话值得拆成两半记。前半句“只拦路”：磁力垂直于速度，所以它改变的是速度箭头的\textbf{指向}。后半句“不助跑”：因为力与每一步的位移都垂直，磁场对电荷\textbf{不做功}，速度箭头的\textbf{长短}不变。想给带电粒子加速，必须请电场出场；磁场的全部本事是当弯道，不是当油门。
\end{mainline}

\step{数学翻译器}

\begin{mainline}
现在把三条定性规律压缩成一个式子。先明确它要回答什么：已知电荷量 \(q\)、速度大小 \(v\)、磁场强弱 \(B\)，以及速度与磁场方向的夹角 \(\theta\)，求磁力有多大。实验给出的答案是
\[
F=qvB\sin\theta .
\]
左边是运动电荷受的磁力；右边每一项各司其职——\(q\) 说“出手认不认你”（不带电就不推），\(v\) 说“你动不动”（不动就不推），\(\sin\theta\) 说“你跑得有多横”（顺着磁感线跑就不推，垂直跑推得最满），\(B\) 则是这片空间“推手的力度”。

这个式子什么时候成立？它是对“点电荷在已知磁场中”这一情形的实验总结，从阴极射线管里的电子到宇宙线里的质子，一百多年的测量都在它的精度范围内，是整个电磁学最可靠的公式之一。什么时候会失效或不够用？后面专门用一节来谈。

先别急着相信它，按本书的规矩，用几个特殊情况当场检查：

\begin{itemize}
  \item \(v=0\)：\(F=0\)。静止电荷不受磁力——开场时挂在磁铁旁边的小球纹丝不动，对上了。
  \item \(\theta=0\) 或 \(\theta=180^\circ\)：\(\sin\theta=0\)，顺着或逆着磁感线飞，磁力为零——先猜一猜的第二题就此判错。
  \item \(\theta=90^\circ\)：\(\sin\theta=1\)，垂直入射时力最大，\(F=qvB\)。极大值出现在“最横”的方向，和直觉实验一致。
  \item \(q\) 换成负数：\(F\) 变号，力的方向整个反转——电子和质子在同一磁场里转弯方向相反，质谱仪就靠这个把它们分开。
  \item \(B\) 加倍：力精确加倍。这是“场强”定义的题中之义：\(B\) 就是按“单位电荷、单位速度受多大横推”标定的，单位是特斯拉。一块强磁铁表面大约 \(0.1\ \mathrm T\) 量级，医院磁共振成像仪的磁场是 \(1.5\) 到 \(3\ \mathrm T\)，地磁场只有约 \(5\times10^{-5}\ \mathrm T\)。
\end{itemize}

方向怎么办？式子只给了大小。方向的规矩是：\(\bm F\)、\(\bm v\)、\(\bm B\) 三者两两垂直，正电荷的指向可以用右手来记——让磁感线穿入手心，四指指向正电荷速度方向，拇指就是力的方向（负电荷则反过来）。用哪只手只是约定，物理内容是“两两垂直”加“正负相反”。

接下来是这个公式最漂亮的推论。磁力处处垂直于速度，于是它只负责拐弯：一个电荷以垂直于匀强磁场的方向射入，每时每刻都被朝侧面推同样大小的力——这正是圆周运动的配置。注意说清楚：这里没有什么新种类的“向心力”，是\textbf{洛伦兹力提供了向心加速度}，角色由磁场扮演，账目仍记在 \(F=ma\) 上。圆周运动要求 \(F=mv^2/r\)，代入 \(F=qvB\)，解出半径
\[
r=\frac{mv}{qB}.
\]
再查一遍特殊情况：速度越快，弯越缓、半径越大——冲得猛的车需要更大的弯，合理；磁场越强，半径越小——推得越狠拐得越急，合理；质量越大半径越大——越重的粒子越“倔”，合理。

更有意思的是转一圈的时间：周长是 \(2\pi r\)，除以速率 \(v\)，得到
\[
T=\frac{2\pi m}{qB}.
\]
速率 \(v\) 竟然从式子里消失了：跑得快的粒子圈子大，跑得慢的圈子小，大家\textbf{同时}回到出发点。这不是巧合，它是一台真实机器的灵魂——回旋加速器就让粒子在磁场里一圈圈转，每过半个圈用电场推一把提速；因为周期不随速度变，电场只要按固定节拍开关，就能一直踩上点。从二十世纪三十年代桌面大小的原型，到今天医院里用来生产放射性药物的小型回旋加速器，用的都是这同一个“节拍免费”的性质。

来做一个带真实数据的估算，体会一下数量级。太阳风里一个质子：质量 \(m=1.67\times10^{-27}\ \mathrm{kg}\)，电荷 \(q=1.6\times10^{-19}\ \mathrm C\)，典型速度 \(v=5\times10^{5}\ \mathrm{m/s}\)。它在地球赤道上空垂直于地磁场（取 \(B=5\times10^{-5}\ \mathrm T\)）飞行时，回旋半径
\[
r=\frac{(1.67\times10^{-27})(5\times10^{5})}{(1.6\times10^{-19})(5\times10^{-5})}
=\frac{8.35\times10^{-22}}{8\times10^{-24}}\approx 1\times10^{2}\ \mathrm m.
\]
一百多米——在地磁场的尺度上，这只算原地打转。也就是说，地球磁场弱归弱，对付太阳风里单个的带电粒子却绰绰有余：它们很难笔直地穿透进来，更多是被迫绕着磁感线打转、顺着线滑走。这就是极光故事的关键一环，稍后就用它。

单电荷的账算清了，还差一步就能解释电动机：把亿万个电荷的横推加起来。电流 \(I\) 的意思是每秒有 \(I\) 库仑的电荷流过导线截面；一段长 \(L\) 的直导线里，电荷淌过这段距离用的电量和速度搭配起来，\(qv\) 恰好可以改写成 \(IL\)。于是整根导线受的磁力是
\[
F=BIL\sin\theta ,
\]
其中 \(\theta\) 是导线（电流方向）与磁场的夹角。老规矩，特殊情况先查一遍：\(I=0\)，不通电就没有力——奥斯特实验里断电后指针乖乖回家；\(\theta=0\)，导线顺着磁场放，力为零；\(\theta=90^\circ\)，导线横架在磁场里，力最大，\(F=BIL\)；电流反向，力的方向整个反转。这个式子有个够分量的应用：配电柜里两根相距 \(d\) 的平行铜排，同向电流会互相吸引，每米长度上的吸力大约是 \(2\times10^{-7}I^{2}/d\) 牛顿。正常运行电流几百安培，这点力无关紧要；可一旦发生短路，电流冲到 \(10^{4}\) 安培量级，两根相距 \(10\) 厘米的铜排之间每米受力约 \(2\times10^{-7}\times10^{8}/0.1=20\) 牛顿，而且全柜的铜排受力叠加、方向一致，支架稍不结实就会被当场掰弯——这就是前面提到“铜排要牢牢锁住”的那笔账。
\end{mainline}

\begin{figure}[htbp]
\centering
\begin{tikzpicture}[scale=1.05]
  % 匀强磁场指向纸内
  \foreach \x in {-2.4,-1.2,0,1.2,2.4}
    \foreach \y in {-2.4,-1.2,0,1.2,2.4}
      \node[blue!55!black] at (\x,\y) {$\otimes$};
  \draw[dashed,thick] (0,0) circle (1.8);
  \fill[red!70!black] (1.8,0) circle (3pt);
  \node[red!70!black,anchor=west] at (1.95,0) {$+q$};
  \draw[-{Stealth[length=3mm]},very thick,green!45!black] (1.8,0) -- (1.8,1.3)
    node[right] {$\bm v$};
  \draw[-{Stealth[length=3mm]},very thick,orange!85!black] (1.8,0) -- (0.6,0)
    node[above] {$\bm F$};
  \fill (0,0) circle (1.2pt);
  \node[anchor=north west] at (-2.5,-2.6) {匀强磁场中垂直入射的正电荷：};
  \node[anchor=north west] at (-2.5,-3.2) {洛伦兹力始终指向圆心，速率不变。};
\end{tikzpicture}
\caption{磁场里的圆轨道：力永远横在速度侧面、指向圆心，只改方向不改快慢。}
\end{figure}

\begin{mathbridge}
用矢量语言，洛伦兹力一行写完：\(\bm F=q\,\bm v\times\bm B\)。叉乘的模长 \(|v||B|\sin\theta\) 正好是主线里的公式，方向由右手定则给出，垂直关系则是叉乘的自带属性——一个算符把三条实验规律全装了进去。

“磁力不做功”也变成一步推导：功率是力与速度的点乘，\(P=\bm F\cdot\bm v=q(\bm v\times\bm B)\cdot\bm v\)。叉乘的结果天然垂直于 \(\bm v\)，垂直向量的点乘恒等于零，所以 \(P\equiv0\)，每一瞬间都为零。这不是近似，是矢量几何的直接判决。

速度与磁场不垂直时，把 \(\bm v\) 拆成两份：沿磁感线的分量 \(v_{\parallel}\) 不受力，保持匀速；垂直分量 \(v_{\perp}\) 照旧画圆。两个运动叠起来，轨迹是一条\textbf{螺旋线}，像拉开的弹簧，一圈圈缠在磁感线上前进。带电粒子沿地磁感线从赤道“滑”向两极，走的就是这种螺旋轨道。范艾伦辐射带里，质子与电子被地磁场这样关了许多年：两极处磁场变强，螺旋的“螺距”被压短直至反向，粒子在两极之间来回弹跳，像被两面磁镜夹着——这是等离子体磁约束的基本画面，人造的托卡马克装置用的也是同一思路，只是磁场形状更讲究。
\end{mathbridge}

\begin{challenge}
磁场还有一个让整个图景倒过来的问题：磁力依赖速度，可速度依赖参考系。你坐在实验室里看一个静止电荷：\(v=0\)，没有磁力。另一个人匀速走过，在他看来同一个电荷明明在动，\(v\neq0\)，应该受磁力。同一个电荷，受力还能因人而异？

这个矛盾没有击垮理论，反而成了狭义相对论最重要的入口之一。严格的分析表明：换参考系时，电场和磁场会互相“掺和”——你认为纯粹的磁场，在运动参考系里会混出一部分电场，而电场恰恰可以作用在静止电荷上。两套账算下来，所有观测者对粒子轨迹的预言完全一致。电与磁不是两种东西，是同一个“电磁场”在不同参考系里呈现的两种侧面。这条线索会在第 24 章麦克斯韦方程那里收紧，并在相对论部分彻底兑现：电磁学内部早已藏着一个要求修改时空观的机关，爱因斯坦只是把它拧开了。
\end{challenge}

\step{模型的边界}

\begin{mainline}
洛伦兹力公式 \(F=qvB\sin\theta\) 的可靠程度罕见地高，但它的使用说明里有几行小字。

第一行：它说的是“一个电荷在给定磁场里受的力”，默认这个电荷自己不会明显搅乱磁场。一个电荷跑起来同样激发磁场，只是那点贡献在绝大多数场合可以忽略。当真要算“两个运动电荷相互的磁力”时，严格做法是分两步——甲的电流产生场、场再推乙——而不是把库仑定律和洛伦兹力直接拼在一起对着两个电荷乱用。

第二行：它只给大小，方向是独立的一条规矩。最常见的误用就是只记 \(qvB\)、忘掉 \(\sin\theta\)，以为“有磁场、有速度就必有磁力”。沿磁感线方向飞的粒子如入无人之境——先猜一猜第二题的错因正在这里。与方向有关的另一条误用是把磁感线当成会交叉的真实丝线：如果两条磁感线在某点相交，那一点的磁场就要同时朝两个方向，小磁针会无所适从——磁场在每一点只有一个确定的方向，所以磁感线永不相交。

第三行：别把“磁力不做功”讲成“磁场什么都不能推动”。电动机明明在输出机械功，发电机明明发出了电，能量从哪儿来？精确的说法是：\textbf{洛伦兹力对单个电荷不做功，但它可以把别的力做的功“转手”}。电动机里，电源把能量交给电荷，磁场像一根拐弯的传动杆，把这份能量转成线圈的转动；磁场自己是中介，不是能源。谁要是做出“只靠永磁体就永远转下去”的机器，他违反的不是某条工程限制，而是这笔能量账本身。发电机的故事更微妙——为什么摇手柄能发电，留到第 23 章电磁感应再揭底。

第四行：公式的经典形象在两头会“不够”。往快里走，粒子速度接近光速时，力公式本身依然成立，但“动量随速度怎么涨”的账要按相对论重算，回旋周期不再是常数——这就是为什么大型加速器要精心设计变频方案或改用别的构型。往小里走，到了原子尺度，磁矩的取值变成量子化的，电子磁矩主要来自内禀自旋而不是轨道转圈，经典图像只能当粗略的脚手架。

最后是一条边界提醒。本章把磁场当作真实存在的场来处理——它能储存能量、能对远处的电荷出手，这已是现代物理的标准做法。但“场究竟是不是一种物质”这类问题，本书不假装有简单答案；我们能肯定的是操作层面的事实：按场的规矩记账，所有预言都对。
\end{mainline}

\step{回头解释开场现象}

\begin{mainline}
回到那台被磁铁揉皱的电视机。显像管里，电子从尾巴的电子枪射出，被电场加速到每秒几万千米，笔直扑向屏幕上的荧光粉点。本来有一套线圈产生的磁场按精确节奏偏转它们，一行一行扫出画面。你把冰箱贴靠近时，等于给电子的航路上空降了一片额外的磁场：每颗电子突然多挨一记横向的洛伦兹力，落点偏移，颜色错位，画面立刻揉成一团。磁铁拿开，额外的力消失，扫描秩序恢复——只有少数老电视因为内部钢件被磁化而留下色斑，那已经是磁化的故事，不是洛伦兹力的现场。

再看极光。太阳风是太阳喷出的带电粒子流，质子和电子以每秒几百千米的速度扑向地球。它们进入地磁场后，无法笔直穿透：垂直于磁感线的运动被掰成半径百来米的回旋，沿磁感线的方向却畅通无阻——磁力管得住“横”，管不住“顺”。于是粒子们只能顺着磁感线螺旋滑行，像弹珠进了漏斗的滑槽。地磁场的磁感线在南北两极附近收拢、扎进地球，粒子也就顺着这些“滑槽”被集中卸到两极上空一百千米左右的高度，撞上氧原子和氮分子，撞出绿色与紫红色的光。极光为什么在两极而不在赤道？因为赤道上空的磁感线平行于地面、横挡着粒子，是两堵墙；极区上空的磁感线竖着插下去，是两扇门。墙与门都是同一个磁场画的。

体检中心的磁共振成像仪、电动交通工具里的电机、实验台上把不同质量粒子分开的质谱仪，都是同一笔账的不同花法：\textbf{只认运动、永远横推、不做功}，三条规矩，撑起一大半现代电磁技术。

顺便回看几个你口袋里的例子。银行卡背面的黑色磁条，是一层可以被逐段磁化的材料，每一段像一枚指向固定的小磁针，刷卡时磁头扫过，读取的就是这些磁矩排出来的图案——这也是磁条卡要远离强磁铁的原因：外磁场会把整排小磁针强行掉头，数据就此抹掉。磁悬浮列车则把磁力用到了另一个极端：用电磁铁之间的吸力或斥力把几百吨的车体托离轨道，消除轮轨摩擦，剩下的工作交给车载电机版的洛伦兹力去完成。从一枚冰箱贴到一列火车，中间没有新物理，只有同一套规矩在不同功率等级上的反复兑现。
\end{mainline}

\step{脑洞实验与挑战题}

\begin{thought}[如果地磁场关一天]
地磁场不是永恒的铁饭碗：岩石记录显示它在地质历史上翻转过许多次，强度也会起伏。假设明天它突然消失一天，我们的生活会怎样？先别想灾难片。指南针集体失业，靠地磁导航的候鸟和信鸽会迷路；近地轨道的卫星失去磁屏蔽，太阳风质子将长驱直入轰击器件，航天器故障率上升。但大气层本身就是一堵厚墙，地表辐射剂量只会小幅上升，我们不会立刻被烤焦。真正有趣的是反向现象：磁场变弱的时候，更多带电粒子漏进高层大气，中低纬度也可能看到极光。这个思想实验的要点在于：一块“弱得可怜”的五万分之一特斯拉的磁场，靠地球级别的体积，撑起了一把全球级的伞——场的强弱，从来要和它占据的空间一起算。
\end{thought}

\begin{puzzles}
  \item 回旋加速器里，粒子每转半圈就被电场推快一次，轨道半径越来越大。有人说：“半径越大，转一圈路越长，周期总会变长，电场节拍迟早踩空。”这个说法对吗？（提示：用 \(r=mv/qB\) 和 \(T=2\pi m/qB\) 检查“路更长”和“时更长”是否同步；再想想这个结论在粒子接近光速时为什么会失效。）
  \item 两根平行导线通以\textbf{相反}方向的电流，会互相排斥。可是把两束电流想象成两队同向奔跑的电荷，直觉上“同方向的伙伴”似乎该相吸。到底哪个才对？（提示：导线是电中性的，跑动的电子之外还有静止的正离子；先用“电流产生磁场、磁场再推另一根导线里的电流”这条两步路线判断方向，再和实验对照。）
  \item 一个带负电的粒子垂直射入指向纸内的匀强磁场，它会朝哪个方向转弯？如果把磁场方向整个翻转，转弯方向变不变？如果把磁场撤掉、换成一个方向合适的匀强电场，能让它走出同样的圆轨道吗？（提示：负电荷受力与右手定则给出的方向相反；圆轨道要求力永远指向圆心，方向恒定的电场做得到吗？）
  \item 扬声器靠磁场推线圈发声，线圈推动纸盆，纸盆推动空气——追到最后，空气确实被推动了。这与“磁力不做功”矛盾吗？（提示：跟踪能量的接力棒：谁把能量交给线圈里的电荷？磁场在接力中扮演的是选手还是交接区？电动机那段的“传动杆”说法可以直接搬过来。）
\end{puzzles}
